\chapter{Historia de la aviación y de la carrera espacial}\label{ch:02}
\objetivos{Relacionar los hitos históricos con los problemas técnicos que resolvieron; distinguir demostración, capacidad operativa y explotación sostenida; analizar fuentes sin convertir una cronología en una explicación causal.}

La historia aeroespacial permite reconocer un patrón: una idea físicamente posible puede tardar décadas en convertirse en un sistema práctico. Para volar se necesitaban conocimientos de fuerzas y materiales, pero también motores suficientemente ligeros, mecanismos de control y una forma organizada de ensayar. Para alcanzar la órbita hicieron falta propulsión, guiado, estructuras y una infraestructura capaz de integrarlos. Este capítulo desarrolla la secuencia presentada en los materiales de la asignatura y la interpreta como historia de problemas de ingeniería \cite{TAE1}.

\section{Antes del vuelo: imaginar, medir y separar funciones}
Los relatos de vuelo expresan un deseo humano; no constituyen evidencia de que sus protagonistas dispusieran de máquinas capaces de cumplirlo. Algo parecido ocurre con los dibujos de Leonardo da Vinci: son relevantes para estudiar imaginación mecánica y observación de las aves, pero un dibujo no equivale a una demostración experimental. La distinción entre propuesta y realización seguirá siendo necesaria al evaluar vehículos contemporáneos.

La imitación directa de un pájaro plantea dificultades de escala, potencia muscular y complejidad mecánica. El paso decisivo hacia el avión consistió en separar funciones. Un elemento produce sustentación, otro proporciona propulsión y otros permiten estabilizar y orientar. En los trabajos de George Cayley aparece esa lectura funcional del vehículo. No hacía falta reproducir todos los movimientos de un ave para satisfacer las mismas leyes físicas.

Los planeadores permitieron estudiar el vuelo sin añadir inmediatamente el problema del motor. Las experiencias de Otto Lilienthal mostraron el valor de repetir ensayos y relacionar forma, incidencia y fuerza aerodinámica, aunque también pusieron de manifiesto los riesgos de un control limitado. Un planeador es un laboratorio en movimiento: su fuente de energía es la disminución de energía potencial, y su duración depende tanto de la resistencia como del entorno atmosférico.

\section{Aerostación y vuelo dirigido}
Los globos de aire caliente de finales del siglo XVIII demostraron que sostener una carga en el aire no exigía necesariamente alas. Su principio era conocido en hidrostática, pero hacerlo operativo requería fabricar una envolvente ligera, producir calentamiento y controlar fugas y cargas. El logro técnico fue una integración de materiales y procedimientos, además de una aplicación de Arquímedes.

El globo libre no podía imponer una trayectoria horizontal respecto al aire. Incorporar una hélice dio lugar a un problema nuevo: conseguir que el empuje superase una resistencia considerable sin añadir demasiada masa. Los dirigibles combinaron gradualmente envolventes, estructuras, motores y superficies de control. Una gran estructura rígida puede conservar la forma y distribuir cargas, pero aumenta complejidad y masa; una envolvente flexible es más sencilla y necesita mantener condiciones internas compatibles con su forma.

Comparar globos, dirigibles y aviones como si pertenecieran a una única competición de velocidad oculta sus propósitos distintos. La permanencia, la carga y la necesidad de infraestructura pueden favorecer soluciones diferentes. La historia no es una escalera en la que toda tecnología nueva elimina a la anterior: cambian los compromisos y los espacios de aplicación.

\section{Los Wright: el vuelo como problema de integración}
La NASA sitúa el vuelo pilotado, controlado y autopropulsado de los hermanos Wright el 17 de diciembre de 1903. Su explicación destaca una secuencia de investigación: revisión de antecedentes, planeadores, ensayos en túnel, mejora del control y desarrollo de la propulsión. El planeador de 1902 incorporó control de los tres ejes; el aparato de 1903 añadió motor y hélices \cite{NASA-BGA}. La formulación completa del hito importa: decir únicamente ``primer objeto que voló'' excluiría sin motivo los aerostatos y los planeadores anteriores.

La enseñanza de ingeniería que extraemos es que la capacidad de sostenerse no basta para realizar una misión. Un aparato que produce más sustentación que su peso, pero no permite controlar la orientación, puede abandonar el suelo y seguir siendo impracticable. La realimentación experimental también es esencial: si una tabla de datos no predice los resultados del prototipo, hay que revisar el ensayo, su escala y la interpretación del coeficiente.

La hélice puede entenderse como un conjunto de alas rotatorias. Ese vínculo explica que mejorar conocimientos aerodinámicos de perfiles beneficie también a la propulsión. Es un ejemplo temprano de transferencia entre subsistemas. El éxito de la integración no elimina la contribución de precursores; muestra que los conocimientos parciales adquieren otro valor al combinarse en un vehículo controlable.

\section{Del prototipo al transporte y al rotor}
Durante las primeras décadas del siglo XX evolucionaron mandos, alas, estructuras y motores. Las estructuras arriostradas ofrecían ligereza con materiales disponibles, a costa de resistencia adicional de montantes y cables. El progreso hacia revestimientos resistentes permitió distribuir cargas de otra forma. Esta evolución no se explica solo por una mejor forma exterior: dependió de materiales, fabricación, cálculo y experiencia de servicio.

Las guerras aceleraron inversión y producción, pero imponer un requisito militar de velocidad o maniobra no equivalía a optimizar transporte civil. Un servicio comercial necesita regularidad, mantenimiento, economía e infraestructura. El desarrollo de la aviación comercial de hélice, cuyo panorama docente ilustra el DC-3, exige por ello estudiar el sistema de explotación además del avión \cite{TAE1}.

El autogiro desarrollado por Juan de la Cierva introdujo soluciones decisivas para un rotor cuyas palas avanzan y retroceden respecto a la corriente. La articulación permite acomodar diferencias de carga. En autorrotación la energía que mantiene el giro procede del flujo a través del rotor; en un helicóptero durante vuelo estacionario ordinario procede del motor. Atribuir la invención de todo el helicóptero a una única persona simplifica un proceso con numerosos desarrollos; Sikorsky fue una figura destacada en la consolidación de una configuración práctica, no el origen de toda idea de ala rotatoria.

\section{Reacción y vuelo rápido}
La propulsión a reacción hizo posible otra combinación de velocidad y altitud. Los desarrollos de Whittle y von Ohain ilustran que un mismo principio físico puede madurar en trabajos independientes. Una turbina de gas integra compresor, combustión y turbina: el reto no es simplemente encender combustible, sino mantener un balance de potencia y un funcionamiento estable con materiales sometidos a temperaturas altas.

Al acercarse a Mach uno aparecen regiones supersónicas y ondas de choque, aunque la velocidad de la aeronave sea todavía subsónica. La llamada barrera del sonido reunió problemas distintos: aumento de resistencia, cambios de momentos, control y cargas. No era una pared material. Comprender la compresibilidad permitió formular soluciones de ala, entrada y propulsión que no se obtenían extrapolando directamente ensayos lentos.

El transporte supersónico mostró que demostrar velocidad no equivale a alcanzar cualquier objetivo económico. Consumo, ruido, mantenimiento y restricciones operativas afectan a la viabilidad. Esta observación es histórica y metodológica: no se deduce que una capacidad carezca de valor porque no se extienda a todos los mercados.

\section{De la astronomía a la órbita}
La observación astronómica produjo modelos cada vez más precisos del movimiento celeste. El heliocentrismo, las leyes de Kepler y la gravitación permitieron describir órbitas con una base matemática. Sin embargo, calcular una trayectoria no proporciona automáticamente un vehículo para recorrerla. La astronáutica necesitó, además, tecnologías de propulsión y control capaces de impartir gran velocidad y resistir el lanzamiento.

La ecuación del cohete vincula incremento de velocidad ideal, velocidad efectiva de escape y razón de masas. Explica por qué expulsar masa funciona en el vacío y por qué la masa estructural es crítica. El empleo de etapas permite abandonar estructuras ya innecesarias. La dificultad de alcanzar órbita no consiste principalmente en subir unos cientos de kilómetros, sino en adquirir una velocidad horizontal elevada sin agotar el vehículo en el proceso.

El contexto de Guerra Fría conectó lanzadores, misiles, prestigio y ciencia. Sputnik 1 fue puesto en órbita el 4 de octubre de 1957; Explorer 1 siguió en 1958 y aportó observaciones asociadas a los cinturones de radiación. NASA comenzó su actividad en octubre de 1958 \cite{NASA-HISTORY}. Estos datos no significan que toda investigación espacial fuese únicamente una aplicación militar: describen un contexto de financiación y competencia con objetivos simultáneos.

\section{Vuelo humano, exploración y continuidad}
El vuelo de Yuri Gagarin en 1961 estableció el primer vuelo espacial humano orbital. Una trayectoria suborbital, como la inicial de Alan Shepard, demuestra capacidades diferentes: supera grandes altitudes sin completar una órbita. Comparar ambos hitos exige definir el criterio, no únicamente contar días entre lanzamientos \cite{TAE1}.

Apollo 11 llevó a Armstrong, Aldrin y Collins a la Luna en julio de 1969; Armstrong y Aldrin descendieron a la superficie y Collins permaneció en órbita lunar \cite{NASA-HISTORY}. El regreso era parte central del objetivo. La misión integró navegación, encuentros orbitales, comunicaciones, soporte vital y protección térmica; un lanzador potente era necesario, pero insuficiente.

Las estaciones orbitales y las sondas ampliaron la permanencia y la investigación. Los satélites de observación, comunicaciones y navegación muestran que la utilidad de una misión no se mide solo por su distancia. También cuentan cobertura, precisión, duración y disponibilidad de los datos. La historia permite reconocer distintas nociones de éxito sin reducirlas a una sola lista de primeros puestos.

\section{Caso resuelto: valorar una propuesta histórica}
\begin{enumerate}
\item \textbf{Enunciado.} Un estudio ficticio compara dos prototipos de comienzos de la aviación. ¿Basta su potencia específica para decidir cuál puede realizar vuelo controlado?
\item \textbf{Datos.} A dispone de $10\,\mathrm{kW}$ y masa total $260\,\mathrm{kg}$; B de $18\,\mathrm{kW}$ y $310\,\mathrm{kg}$. Solo A tiene ensayos documentados de mandos en tres ejes.
\item \textbf{Incógnitas.} Potencia por masa y conclusión técnica justificable.
\item \textbf{Esquema verbal.} Se comparan capacidad energética e información de control, manteniendo separadas las dos preguntas.
\item \textbf{Hipótesis.} Masas completas comparables y potencias al eje medidas del mismo modo; el ejemplo no representa un aparato histórico concreto.
\item \textbf{Modelo.} Indicador $P/m$ complementado con una evaluación cualitativa de evidencia.
\item \textbf{Ecuación.} $\pi_P=P/m$, cuya unidad es $\mathrm{W\,kg^{-1}}$; no es un coeficiente adimensional.
\item \textbf{Despeje.} Se divide la potencia de cada prototipo por su propia masa completa.
\item \textbf{Sustitución.} $\pi_{P,A}=10000/260$ y $\pi_{P,B}=18000/310$.
\item \textbf{Resultado.} A ofrece $38.5\,\mathrm{W\,kg^{-1}}$ y B $58.1\,\mathrm{W\,kg^{-1}}$; B supera a A en este indicador.
\item \textbf{Interpretación.} No se puede concluir que B realice vuelo controlado. Falta evidencia de sustentación, resistencia, transmisión, estabilidad y autoridad de los mandos.
\item \textbf{Comprobación independiente.} La razón $(18/10)(260/310)=1.51$ coincide con $\pi_{P,B}/\pi_{P,A}$. La comparación cualitativa se verifica preguntando qué afirmación permite realmente cada documento.
\item \textbf{Limitaciones.} Sin mapas de hélice ni datos aerodinámicos, tampoco se puede afirmar que A vuele. Haber ensayado los mandos reduce una incertidumbre y no resuelve todas las demás.
\end{enumerate}

\section{Ejercicios propuestos}
\begin{enumerate}
\item \textbf{Básico.} Ordenar planeador, vuelo propulsado controlado y transporte regular por la capacidad adicional que requiere cada transición.
\item \textbf{Intermedio.} Explicar qué información faltaría para comparar dos motores históricos a partir de su masa en seco.
\item \textbf{Examen.} Construir un argumento de ocho líneas sobre por qué llegar al espacio y permanecer en órbita son objetivos diferentes. Incluir una condición física y una consecuencia de diseño.
\item \textbf{Conceptual.} Una fuente atribuye un ``primer vuelo'' sin explicar si estaba tripulado ni si era controlado. ¿Qué comprobaciones documentales haría?
\end{enumerate}
\section{Respuestas y orientación}
1. El motor añade energía sostenida; el servicio regular añade fiabilidad, mantenimiento y operación. 2. Potencia útil, combustible, instalación y duración de funcionamiento, entre otros. 3. La órbita requiere velocidad tangencial compatible con caída libre; la trayectoria suborbital intersecta de nuevo regiones bajas. 4. Precisar definición del hito, fecha, testigos, registros y contexto, y contrastar fuentes.

\resumen{La evolución aeroespacial combina ideas físicas, mediciones, materiales, control e infraestructura. Las fechas ubican hitos; la explicación técnica identifica qué problema se resolvió y qué seguía pendiente. Una demostración no garantiza capacidad operativa, y un indicador favorable no sustituye la integración del vehículo.}
\section{Autoevaluación}
\begin{enumerate}
\item ¿Por qué separar sustentación y propulsión fue una idea fértil?
\item ¿Qué enseñan los Wright sobre el valor de revisar datos experimentales?
\item ¿Cómo cambia una comparación histórica al definir exactamente el éxito?
\end{enumerate}
